\documentclass[a4paper]{article}
\usepackage{amsfonts}
\usepackage{amsmath}
\usepackage{amssymb}
\usepackage{graphicx}

\usepackage{cancel}

\begin{document}
  
\noindent \large Auteur: Siebe Mees \\
\noindent \large Studierichting: Informatica\\
\noindent \large Oefeningenreeks 6G nr. 2.8\\

\medskip

\normalsize

$f(x)=e^{|x|}$\\

$f(-x)=e^{|-x|}=e^{|x|}=f(x)$\\

$\Rightarrow f $ is even $ \Rightarrow b_n=0$\\

$f(x)\sim \dfrac{a_0}{2}+\displaystyle\sum_{n=1}^{\infty}(a_n\cos(nx))$\\

$a_0=\dfrac{1}{\pi}\displaystyle\int_{-\pi}^{\pi}f(x)dx=\dfrac{1}{\pi}\displaystyle\int_{-\pi}^{\pi}e^{|x|}dx$\\

Omdat $f$ even is:\\

$ = 2\cdot\dfrac{1}{\pi}\displaystyle\int_{0}^{\pi}e^xdx$\\

$=\dfrac{2}{\pi}\left[e^x\right]_{0}^{\pi}$\\

$ \Rightarrow a_0 =\dfrac{2}{\pi}(e^{\pi}-1)$\\

$a_n=\dfrac{1}{\pi}\displaystyle\int_{-\pi}^{\pi}f(x)\cos(nx)dx=\dfrac{1}{\pi}\displaystyle\int_{-\pi}^{\pi}e^{|x|}\cos(nx)dx$\\

Omdat $f$ even is:\\

$=\dfrac{2}{\pi}\displaystyle\int_{0}^{\pi}e^x\cos(nx)dx$\\

$ \left \{
\begin{array}{l}
u=\cos(nx) \\
dv=e^xdx
\end{array}
\right .
\Rightarrow
\left \{
\begin{array}{l}
du=-n\sin(nx)dx \\
v=e^x
\end{array}
\right . $\\

$I=e^x\cos(nx)-n\displaystyle\int_{0}^{\pi} e^x\sin(nx)dx$\\

$ \left \{
\begin{array}{l}
u=\sin(nx) \\
dv=e^xdx
\end{array}
\right .
\Rightarrow
\left \{
\begin{array}{l}
du=n\cos(nx)dx \\
v=e^x
\end{array}
\right . $\\

$J=e^x\sin(nx)-n\displaystyle\int_{0}^{\pi} e^x\cos(nx)dx$\\

$\Rightarrow I=e^x\cos(nx)-n(e^x\sin(nx)-nI)$\\

$=e^x\cos(nx)-ne^x\sin(nx)+n^2I$\\

$I(1+n^2)=e^x\cos(nx)-ne^x\sin(nx)$\\

$=e^x(\cos(nx)+n\sin(nx))$\\

$ = \dfrac{2}{\pi}\left[\dfrac{e^x(\cos(nx)+n\sin(nx))}{1+n^2}\right]_{0}^{\pi}$\\

$=\dfrac{2}{\pi}\left(\dfrac{e^{\pi}(-1)^n}{1+n^2}-\dfrac{1}{1+n^2}\right)$\\

$\Rightarrow a_n =\dfrac{2(e^{\pi}(-1)^n-1)}{\pi(1+n^2)}$\\

$f(x)\sim \dfrac{e^{\pi}-1}{\pi}+\displaystyle\sum_{n=1}^{\infty}\dfrac{2(e^{\pi}(-1)^n-1)}{\pi(1+n^2)}\cos(nx)$\\

\begin{thebibliography}{999}
%\bibitem{a} Referentie
\end{thebibliography}
\end{document} 
